\section*{Chapitre \ref{ch:b2:reduction} --- Réduction des endomorphismes}

\begin{solution}{exo:b2:reduction:1}
$A$~: $\chi_A = X^2 - 2X - 3 = (X - 3)(X + 1)$. Vecteurs propres~: pour
$3$~: $(1,1)$~; pour $-1$~: $(1,-1)$. Donc $P = \begin{pmatrix} 1 & 1\\ 1
& -1\end{pmatrix}$ donne $P^{-1}AP = \operatorname{diag}(3, -1)$.

$B = J - I$ où $J$ est la matrice remplie de $1$. $J$ est de rang $1$
avec $Jv = 3v$ pour $v = (1,1,1)$ et $Jw = 0$ sur le plan $x + y + z =
0$~: le \omterm{def:b2:reduction:eigen}{spectre} de $B$ est $\{2, -1\}$ avec sous-espaces propres
$\operatorname{Vect}(1,1,1)$ (dimension $1$) et $\{x + y + z = 0\}$
(dimension $2$, base $(1,-1,0), (1,0,-1)$). $P$ ayant ces trois
colonnes donne $P^{-1}BP = \operatorname{diag}(2, -1, -1)$.
\end{solution}

\begin{solution}{exo:b2:reduction:2}
\emph{Sous-espaces propres~:} $\chi_C = (X-1)^2$, \omterm{def:b2:reduction:eigen}{valeur propre} unique
$1$~; $\ker(C - I) = \ker\begin{pmatrix} 0&1\\ 0&0\end{pmatrix}$ est la
droite $\operatorname{Vect}(e_1)$~: dimension $1 < 2 = m_1$, donc non
\omterm{def:b2:reduction:diag}{diagonalisable} (\cref{thm:b2:reduction:diagcrit}).

\emph{\omterm{def:b2:reduction:polyu}{Polynôme minimal}~:} $\mu_C$ divise $(X-1)^2$ et $C \neq I$, donc
$\mu_C = (X-1)^2$~: une racine double, donc non \omterm{def:b2:reduction:diag}{diagonalisable}
(\cref{cor:b2:reduction:minpolycrit}).
\end{solution}

\begin{solution}{exo:b2:reduction:3}
$X^2 - 5X + 6 = (X-2)(X-3)$~: scindé à racines simples, donc $u$ est
\omterm{def:b2:reduction:diag}{diagonalisable} (\cref{cor:b2:reduction:minpolycrit}), avec
$\operatorname{Sp}(u) \subseteq \{2, 3\}$. \omterm{def:b2:reduction:eigen}{Spectres} possibles~:
$\{2\}$ ($u = 2\,\mathrm{id}$), $\{3\}$ ($u = 3\,\mathrm{id}$), ou
$\{2, 3\}$.

Puissances~: on cherche $u^k = a_k\,\mathrm{id} + b_k\,u$. Sur les
sous-espaces propres, cela s'écrit $2^k = a_k + 2b_k$ et $3^k = a_k +
3b_k$~: en résolvant, $b_k = 3^k - 2^k$, $a_k = 3\cdot2^k - 2\cdot
3^k$~:
\[
u^k = (3\cdot 2^k - 2\cdot 3^k)\,\mathrm{id} + (3^k - 2^k)\, u .
\]
(Valable pour les trois \omterm{def:b2:reduction:eigen}{spectres}~: les identités valent \omterm{def:b2:reduction:eigen}{valeur propre}
par \omterm{def:b2:reduction:eigen}{valeur propre}.)
\end{solution}

\begin{solution}{exo:b2:reduction:4}
$u$ \omterm{def:b2:reduction:diag}{diagonalisable}~: $P = \prod_{\lambda}(X - \lambda)$ sur le \omterm{def:b2:reduction:eigen}{spectre}
annule $u$, est scindé, à racines simples. Alors $P(u|_F) = P(u)|_F =
0$~: la restriction est annulée par un polynôme scindé à racines
simples, donc \omterm{def:b2:reduction:diag}{diagonalisable}
(\cref{cor:b2:reduction:minpolycrit}).
\end{solution}

\begin{solution}{exo:b2:reduction:5}
$\chi_A = X^2 - X - 1$, racines $\varphi$ et $\psi$ (distinctes)~:
\omterm{def:b2:reduction:diag}{diagonalisable}, de vecteurs propres $(\varphi, 1)$ et $(\psi, 1)$. La
récurrence donne $\begin{pmatrix} F_{n+1}\\ F_n \end{pmatrix} =
A^n \begin{pmatrix}1\\ 0\end{pmatrix}$. Décomposons $(1, 0)$ sur les
vecteurs propres~: $(1,0) = \frac{1}{\varphi - \psi}\bigl((\varphi, 1) -
(\psi, 1)\bigr)$ avec $\varphi - \psi = \sqrt5$. Appliquer $A^n$
multiplie chaque composante propre par la puissance $n$-ième de sa
\omterm{def:b2:reduction:eigen}{valeur propre}~; en lisant la seconde coordonnée~:
\[
F_n = \frac{\varphi^n - \psi^n}{\sqrt 5} .
\]
(Vérification~: $n = 1$ donne $\frac{\varphi - \psi}{\sqrt5} = 1$.)
\end{solution}

\begin{solution}{exo:b2:reduction:6}
Soit $P = \prod_i (X - \mu_i)$ annulant $u^2$, scindé à racines
simples $\mu_i$ (le \omterm{def:b2:reduction:eigen}{spectre} de $u^2$). Puisque $u$ est inversible, $0$
n'est pas \omterm{def:b2:reduction:eigen}{valeur propre} de $u^2$ ($\det u^2 = (\det u)^2 \neq 0$), donc
tous les $\mu_i \neq 0$. Alors
\[
Q(X) = \prod_i (X^2 - \mu_i) = \prod_i (X - \sqrt{\mu_i})(X +
\sqrt{\mu_i})
\]
annule $u$~: $\;Q(u) = \prod_i (u^2 - \mu_i\,\mathrm{id}) =
P(u^2) = 0$. Ses racines $\pm
\sqrt{\mu_i}$ (racines carrées complexes) sont deux à deux distinctes
car les $\mu_i$ sont distincts et non nuls ($\sqrt{\mu_i} =
-\sqrt{\mu_j}$ donnerait $\mu_i = \mu_j$). Scindé + racines simples~:
$u$ est \omterm{def:b2:reduction:diag}{diagonalisable}.

Contre-exemple sans inversibilité~: $u = \begin{pmatrix} 0 & 1\\
0 & 0\end{pmatrix}$~: $u^2 = 0$ est \omterm{def:b2:reduction:diag}{diagonalisable}, $u$ ne l'est pas.
\end{solution}

\begin{solution}{exo:b2:reduction:7}
$A = 2I + N$ avec $N$ le décalage ($N e_2 = e_1$, $Ne_3 = e_2$),
$N^3 = 0$, $N^2 = E_{13}$~: c'\emph{est} la \omterm{thm:b2:reduction:dunford}{décomposition de Dunford}
($2I$ diagonale, $N$ nilpotente, elles commutent~; l'unicité en fait
la bonne). Binôme à termes qui commutent~:
\[
A^k = 2^k I + k 2^{k-1} N + \binom k2 2^{k-2} N^2
= \begin{pmatrix}
2^k & k2^{k-1} & \binom k2 2^{k-2}\\
0 & 2^k & k2^{k-1}\\
0 & 0 & 2^k
\end{pmatrix},
\]
\[
\eu^{tA} = \eu^{2t}\Bigl(I + tN + \frac{t^2}{2}N^2\Bigr)
= \eu^{2t}\begin{pmatrix}
1 & t & t^2/2\\
0 & 1 & t\\
0 & 0 & 1
\end{pmatrix}.
\]
\end{solution}

\begin{solution}{exo:b2:reduction:8}
$X^k - 1$ annule $A$ et est scindé sur $\C$ avec les $k$ racines
distinctes $\eu^{2\iu\pi j/k}$~: $A$ est \omterm{def:b2:reduction:diag}{diagonalisable}
(\cref{cor:b2:reduction:minpolycrit}) et ses \omterm{def:b2:reduction:eigen}{valeurs propres}, racines
de $X^k - 1$, sont des racines $k$-ièmes de l'unité.

Si de plus $A$ est semblable à une matrice triangulaire à diagonale
unité~: toutes les \omterm{def:b2:reduction:eigen}{valeurs propres} valent $1$, et $A$, \omterm{def:b2:reduction:diag}{diagonalisable}
de seule \omterm{def:b2:reduction:eigen}{valeur propre} $1$, vaut $P\,I\,P^{-1} = I$.
\end{solution}

\begin{solution}{exo:b2:reduction:9}
Écrivons $E = \bigoplus_\lambda E_\lambda(u)$
(\cref{thm:b2:reduction:diagcrit}). Chaque $E_\lambda(u)$ est
$v$-stable~: pour $x \in E_\lambda$, $u(v(x)) = v(u(x)) = \lambda
v(x)$. La restriction de $v$ à $E_\lambda(u)$ est \omterm{def:b2:reduction:diag}{diagonalisable}
(\cref{exo:b2:reduction:4})~: on choisit une base de $E_\lambda(u)$
faite de vecteurs propres de $v$. En concaténant ces bases sur tous
les $\lambda$, on obtient une base de $E$ dont les vecteurs sont
vecteurs propres \emph{à la fois} de $u$ (par appartenance à
$E_\lambda(u)$) et de $v$ (par construction).
\end{solution}

\begin{solution}{exo:b2:reduction:10}
($\Rightarrow$) Soit $u$ \omterm{def:b2:reduction:diag}{diagonalisable} et $F$ stable. Alors $u|_F$
est \omterm{def:b2:reduction:diag}{diagonalisable} (\cref{exo:b2:reduction:4})~: $F$ a une base de
vecteurs propres, qui se prolonge, à l'intérieur de chaque sous-espace
propre global $E_\lambda$, en une base de $E_\lambda$ (théorème de la
base incomplète dans $E_\lambda$, en partant de la partie de la base
de $F$ qui s'y trouve --- noter $F = \bigoplus_\lambda (F \cap
E_\lambda)$ puisque $u|_F$ est \omterm{def:b2:reduction:diag}{diagonalisable}). Les vecteurs ajoutés
engendrent un supplémentaire stable (chacun est dans un certain
$E_\lambda$, donc leur engendré est $u$-stable).

($\Leftarrow$) Soit $F = \sum_\lambda E_\lambda(u)$ (un sous-espace
stable) et $G$ un supplémentaire stable. Si $G \neq \{0\}$~:
$\chi_{u|_G}$ est scindé sur $\C$, donc $u|_G$ a un vecteur propre $x
\in G$ (\cref{thm:b2:reduction:trigonalization} ou directement
l'existence d'une racine)~; mais tout vecteur propre de $u$ est dans
$F$, donc $x \in F \cap G = \{0\}$~: contradiction. Ainsi $G = \{0\}$
et $E = F$~: les sous-espaces propres remplissent $E$, c'est-à-dire
$u$ est \omterm{def:b2:reduction:diag}{diagonalisable}.
\end{solution}

\begin{solution}{exo:b2:reduction:11}
Trigonaliser~: $A = PTP^{-1}$, $T = \begin{pmatrix} \lambda & c\\ 0 &
\mu\end{pmatrix}$, $\abs\lambda, \abs\mu < 1$. Alors $A^k =
PT^kP^{-1}$, et il suffit que $T^k \to 0$.

\emph{\omterm{def:b2:reduction:eigen}{Valeurs propres} distinctes~:} la récurrence donne
\[
T^k = \begin{pmatrix}
\lambda^k & c\,\dfrac{\lambda^k - \mu^k}{\lambda - \mu}\\[4pt]
0 & \mu^k
\end{pmatrix},
\]
et chaque coefficient tend vers $0$ ($\abs{\lambda}^k, \abs\mu^k \to
0$).

\emph{\omterm{def:b2:reduction:eigen}{Valeurs propres} égales ($\mu = \lambda$)~:} $T = \lambda I +
cE_{12}$ et $T^k = \lambda^k I + k\lambda^{k-1}cE_{12}$~; le
coefficient $k\lambda^{k-1} \to 0$ puisque $\abs\lambda < 1$ (la
géométrique l'emporte sur le polynôme). Dans les deux cas $T^k \to 0$
coefficient par coefficient, d'où $A^k = PT^kP^{-1} \to 0$ (la
multiplication matricielle par $P, P^{-1}$ fixés est continue dans les
coefficients --- chaque coefficient du produit est une combinaison
linéaire fixée).
\end{solution}

\begin{solution}{exo:b2:reduction:12}
$\ker u$ est de dimension $n - 1$ (théorème du rang), donc $0$ est une
\omterm{def:b2:reduction:eigen}{valeur propre} de \omterm{def:b2:reduction:charpoly}{multiplicité géométrique} $n - 1$, et $\chi_u$ est
divisible par $X^{n-1}$
(\cref{def:b2:reduction:charpoly}~: géométrique $\leq$ algébrique). En
écrivant $\chi_u = X^{n-1}(X - \alpha)$~; le coefficient de $X^{n-1}$
valant $-\operatorname{tr} u$, on obtient $\alpha =
\operatorname{tr} u$~: $\chi_u = X^{n-1}(X - \operatorname{tr} u)$.

Si $\operatorname{tr} u \neq 0$~: la \omterm{def:b2:reduction:eigen}{valeur propre} $\operatorname{tr}
u$ est une racine de $\chi_u$, donc elle porte un vecteur propre~; les
sous-espaces propres pour $0$ et $\operatorname{tr} u$ ont pour
dimensions $n - 1$ et $\geq 1$, de somme $\geq n$~: ils remplissent
$E$, et $u$ est \omterm{def:b2:reduction:diag}{diagonalisable}
(\cref{thm:b2:reduction:diagcrit}). Si $\operatorname{tr} u = 0$~: par
l'\cref{exo:b2:linalg:5}, $u^2 = (\operatorname{tr} u)u = 0$ avec $u
\neq 0$~: $u$ est un nilpotent non nul, et un nilpotent \omterm{def:b2:reduction:diag}{diagonalisable}
est nul (\cref{prop:b2:reduction:nilpotent})~: non \omterm{def:b2:reduction:diag}{diagonalisable}.
\end{solution}

\begin{solution}{pb:b2:reduction:1}
\textbf{1.} Les $k - 1$ premières coordonnées de $Cv_n$ sont
$u_{n+1}, \dots, u_{n+k-1}$ (la surdiagonale décale), et la dernière
est $a_0 u_n + \dots + a_{k-1}u_{n+k-1}$. Donc $v_{n+1} = Cv_n$ vaut
pour tout $n$ si et seulement si les dernières coordonnées coïncident
pour tout $n$, c'est-à-dire si et seulement si $(\mathcal R)$ vaut. En
itérant, $v_n = C^nv_0$.

\textbf{2.} Développons $D_k(X) = \det(XI_k - C)$ selon la première
colonne~: les deux entrées non nulles sont $X$ (position $(1,1)$) et
$-a_0$ (position $(k,1)$). Le premier mineur est de forme $D_{k-1}$
pour les coefficients $a_1, \dots, a_{k-1}$~; le second mineur est
triangulaire supérieur de diagonale $-1$~: \omterm{def:b2:linalg:det}{déterminant} $(-1)^{k-1}$,
de signe $(-1)^{k+1}$ dû à la position. La récurrence sur $k$ (base
$k = 1$~: $X - a_0$) donne
\[
D_k(X) = X\bigl(X^{k-1} - a_{k-1}X^{k-2} - \dots - a_1\bigr) -
a_0 = P(X).
\]
Pour $\mu_C$~: puisque $Q(C^{\mathsf T}) = Q(C)^{\mathsf T}$ pour tout
polynôme, $C$ et $C^{\mathsf T}$ ont les mêmes \omterm{def:b2:linalg:annihilator}{annulateurs}, donc le
même \omterm{def:b2:reduction:polyu}{polynôme minimal}. Pour $C^{\mathsf T}$~: les colonnes donnent
$C^{\mathsf T}e_1 = e_2$, \dots, $C^{\mathsf T}e_{k-1} = e_k$, donc
$(e_1, C^{\mathsf T}e_1, \dots, (C^{\mathsf T})^{k-1}e_1)$ est la base
canonique~: libre. Un polynôme $Q \neq 0$ de degré $< k$ a alors
$Q(C^{\mathsf T})e_1 \neq 0$ (c'est une combinaison non triviale de
vecteurs de base)~: $\deg\mu \geq k$. Comme $\mu \mid \chi = P$ avec
$\deg P = k$~: $\mu_C = P$.

\textbf{3.} Pour $v = (1, \lambda, \dots, \lambda^{k-1})^{\mathsf
T}$~: les lignes $1$ à $k-1$ de $Cv$ donnent $\lambda, \lambda^2,
\dots, \lambda^{k-1}$, c'est-à-dire $\lambda$ fois les $k - 1$
premières entrées de $v$~; la dernière ligne donne $\sum_m
a_m\lambda^m = \lambda^k - P(\lambda) = \lambda^k =
\lambda\cdot\lambda^{k-1}$. Donc $Cv = \lambda v$. Réciproquement, les
équations $(Cx)_i = \lambda x_i$ pour $i < k$ donnent $x_{i+1} =
\lambda x_i$~: tout vecteur propre est proportionnel à $v$ --- tout
sous-espace propre est de dimension exactement $1$. \omterm{def:b2:reduction:diag}{Diagonalisable} si
et seulement si les dimensions des sous-espaces propres somment à $k$
(\cref{thm:b2:reduction:diagcrit}) si et seulement s'il y a $k$
\omterm{def:b2:reduction:eigen}{valeurs propres} distinctes si et seulement si $P$ a $k$ racines
distinctes (les \omterm{def:b2:reduction:eigen}{valeurs propres} sont les racines de $\chi_C = P$).

\textbf{4.} Chaque $(\lambda_i^n)_n$ résout $(\mathcal R)$~:
$\lambda_i^{n+k} = \lambda_i^n\,\lambda_i^k =
\lambda_i^n\sum_m a_m\lambda_i^m$. Liberté~: une combinaison nulle
$\sum_i c_i\lambda_i^n = 0$ pour $n = 0, \dots, k-1$ est un système de
Vandermonde (\cref{exo:b2:linalg:11}) en les $c_i$~: tous $c_i = 0$.
L'espace des solutions est de dimension $k$ (question 6, dont la
preuve est élémentaire et indépendante)~: $k$ solutions libres
forment une base, et les coordonnées sont uniques.

\textbf{5.} $P = X^2 - X - 6 = (X - 3)(X + 2)$~: solution générale
$u_n = A\,3^n + B(-2)^n$. Conditions initiales~: $A + B = 1$, $3A -
2B = 8$~: $A = 2$, $B = -1$~:
\[
u_n = 2\cdot 3^n - (-2)^n .
\]
(Vérification~: $u_2 = u_1 + 6u_0 = 14$ et $2\cdot9 - 4 = 14$.)

\textbf{6.} $P(S)\bigl((u_n)\bigr)$ est la suite $n \mapsto u_{n+k} -
a_{k-1}u_{n+k-1} - \dots - a_0u_n$~: elle s'annule si et seulement si
$(\mathcal R)$ vaut, donc l'ensemble des solutions est $\ker P(S)$, un
sous-espace. L'application $\ker P(S) \to \C^k$, $u \mapsto (u_0,
\dots, u_{k-1})$, est linéaire, injective (la récurrence détermine
$u_{k}, u_{k+1}, \dots$ à partir des $k$ premières valeurs, par
récurrence) et surjective (on définit $u_n$ récursivement à partir de
n'importe quelle donnée initiale)~: dimension $k$.

\textbf{7.} La preuve du~\cref{thm:b2:reduction:kernels} n'utilise
que~: l'identité de Bézout dans $\C[X]$, et le fait que les polynômes
en un endomorphisme fixé commutent. Ni l'une ni l'autre ne mentionne
la dimension de l'espace ambiant~: le lemme vaut mot pour mot pour $S
\in \mathcal{L}(\mathcal{S})$. Ainsi
\[
\ker P(S) = \bigoplus_{i=1}^{r}
\ker\,(S - \lambda_i\,\mathrm{id})^{m_i}.
\]

\textbf{8.} Pour $Q \in \C[X]$~: $(S -
\lambda)\bigl(Q(n)\lambda^n\bigr)_n$ a pour terme $n$-ième
$Q(n{+}1)\lambda^{n+1} - \lambda Q(n)\lambda^n =
\lambda^{n+1}(\Delta Q)(n)$, avec $\Delta Q = Q(X{+}1) - Q(X)$ de
degré $\deg Q - 1$ (les termes de tête se compensent). En itérant,
$(S - \lambda)^m\bigl(Q(n)\lambda^n\bigr) =
\bigl(\lambda^{n+m}(\Delta^m Q)(n)\bigr)_n$, et $\Delta^m Q = 0$
lorsque $\deg Q \leq m - 1$~: l'ensemble de droite est contenu dans le
noyau. C'est un sous-espace de dimension $m$~: les suites
$(n^j\lambda^n)_n$, $0 \leq j < m$, sont libres, car
$\sum_j c_j n^j\lambda^n = 0$ pour tout $n$ force (en divisant par
$\lambda^n \neq 0$) le polynôme $\sum_j c_jX^j$ à s'annuler en tout $n
\in \N$, donc à être nul. Réciproquement $\dim\ker(S - \lambda)^m \leq
m$~: en développant $(S - \lambda)^m = \sum_j
\binom mj(-\lambda)^{m-j}S^j$, l'équation $(S - \lambda)^m u = 0$ est
une récurrence linéaire d'\omterm{def:b2:structures:generated}{ordre} $m$ (coefficient de tête $1$), donc
$u$ est déterminée par $u_0, \dots, u_{m-1}$ comme à la question 6.
L'égalité des dimensions conclut.

\textbf{9.} On combine les questions 7 et 8~: toute solution se
décompose de manière unique en somme d'éléments des $\ker(S -
\lambda_i)^{m_i}$, c'est-à-dire $u_n = \sum_i Q_i(n)\lambda_i^n$ avec
$\deg Q_i \leq m_i - 1$~; les $Q_i$ sont uniques car la décomposition
est directe et, dans chaque facteur, les coefficients de $Q_i$ sont
des coordonnées dans la base $(n^j\lambda_i^n)_j$ (question 8).
Vérification de cohérence sur les dimensions~: $\sum_i m_i = k$.

\textbf{10.} $P = X^2 - 4X + 4 = (X - 2)^2$~: solutions $(a +
bn)2^n$. Données initiales~: $a = 1$, $2(a + b) = 0$, donc $b = -1$~:
\[
u_n = (1 - n)\,2^n .
\]
Vérification~: $u_2 = 4u_1 - 4u_0 = -4$, et $(1 - 2)\cdot4 = -4$.

\textbf{11.} Écrivons $u_n = \lambda_1^n\bigl(c_1 + \sum_{i\geq2}
c_i(\lambda_i/\lambda_1)^n\bigr)$~; chaque rapport est de module $<
1$, donc le crochet tend vers $c_1 \neq 0$~: $u_n \sim c_1\lambda_1^n$.
En particulier $u_n \neq 0$ pour $n$ grand, et
\[
\frac{u_{n+1}}{u_n} =
\lambda_1\,\frac{c_1 + o(1)}{c_1 + o(1)} \longrightarrow
\lambda_1 .
\]

\textbf{12.} On calcule~:
\[
q(a_{n+1}, b_{n+1}) = (a_n + 2b_n)^2 - 2(a_n + b_n)^2
= -a_n^2 + 2b_n^2 = -q(a_n, b_n).
\]
Avec $q(a_0, b_0) = 1 - 2 = -1$~: $a_n^2 - 2b_n^2 = (-1)^{n+1}$.
Structurellement~: $q(a, b) = (a - \sqrt2\,b)(a + \sqrt2\,b)$ et
l'application linéaire $M$ multiplie le facteur $a + \sqrt2 b$ par $1 +
\sqrt2$ et le facteur $a - \sqrt2 b$ par $1 - \sqrt2$ (calculer~:
$a_{n+1} + \sqrt2 b_{n+1} = (1 + \sqrt2)(a_n + \sqrt2 b_n)$)~; le
produit est multiplié par $(1 + \sqrt2)(1 - \sqrt2) = -1 = \det M$ à
chaque pas.

\textbf{13.} Puisque $a_n^2 - 2b_n^2 = (a_n - \sqrt2 b_n)(a_n +
\sqrt2 b_n) = (-1)^{n+1}$,
\[
\Bigl|\frac{a_n}{b_n} - \sqrt2\Bigr|
= \frac{\abs{a_n^2 - 2b_n^2}}{b_n(a_n + \sqrt2 b_n)}
= \frac{1}{b_n(a_n + \sqrt2 b_n)} \leq \frac1{2b_n^2},
\]
en utilisant $a_n \geq b_n \geq 1$ (par récurrence~: les deux
croissent) donc $a_n + \sqrt2 b_n \geq (1 + \sqrt2)b_n \geq 2b_n$.
\omterm{def:b2:reduction:eigen}{Valeurs propres} de $M$~: $\chi_M = X^2 - 2X - 1$, racines $1 \pm
\sqrt2$~; puisque $(a_0, b_0)$ a une composante non nulle sur le
vecteur propre dominant (toutes les entrées positives), $b_n \sim c(1
+ \sqrt2)^n$ avec $c > 0$ (question 11). D'où l'erreur est $\asymp (1
+ \sqrt2)^{-2n} = (3 + 2\sqrt2)^{-n}$~: décroissance géométrique de
rapport $1/(3 + 2\sqrt2) = 3 - 2\sqrt2 \approx 0{,}172$.

\textbf{14.} (a) À partir de la question 9~: $\abs{u_n} \leq \sum_i
\abs{Q_i(n)}\abs{\lambda_i}^n \leq \bigl(\sum_i
\abs{Q_i(n)}\bigr)\rho^n$, et chaque $\abs{Q_i(n)} \leq C_i
n^{m_i - 1} \leq C_i n^{m-1}$ pour $n \geq 1$~: on somme les
constantes. (b) Soit $\rho' = \max_{i \geq 2}\abs{\lambda_i} <
\abs{\lambda_1}$ et $d = \deg Q_1$, de coefficient de tête $c \neq 0$.
Alors $u_n = Q_1(n)\lambda_1^n + R_n$ avec $\abs{R_n} \leq
Cn^{m-1}\rho'^n$, et
\[
\frac{R_n}{Q_1(n)\lambda_1^n} = O\Bigl(n^{m-1-d}
\bigl(\rho'/\abs{\lambda_1}\bigr)^n\Bigr) \longrightarrow 0
\]
(la géométrique l'emporte sur le polynôme). Donc
\[
u_n \sim Q_1(n)\,\lambda_1^n \sim c\,n^d\lambda_1^n,
\qquad
\frac{u_{n+1}}{u_n} \longrightarrow \lambda_1
\quad\text{(puisque } Q_1(n{+}1)/Q_1(n) \to 1\text{)}.
\]
Vérification sur la question 10~: pour $u_n = (1-n)2^n$ le rapport est
\[
\frac{(-n)2^{n+1}}{(1-n)2^n} = 2\,\frac{-n}{1-n}
\longrightarrow 2 = \lambda_1 .
\]

\textbf{15.} Récurrence sur $n$. Pour $n = 1$, $A_{ij}$ compte les
chemins de longueur $1$. Étape~: un chemin de longueur $n + 1$ de $i$
à $j$ est un chemin de longueur $n$ de $i$ vers un sommet $\ell$ suivi
d'une arête $\ell j$~:
\[
\#\{\text{chemins}\} = \sum_{\ell} (A^n)_{i\ell}A_{\ell j} =
(A^{n+1})_{ij}.
\]

\textbf{16.} $J = 3\Pi$ où $\Pi = J/3$ est la projection sur
$\operatorname{Vect}(1,1,1)$ parallèlement au plan $x + y + z = 0$
($\Pi^2 = \Pi$ puisque $J^2 = 3J$). Alors $A = J - I = 2\Pi - (I -
\Pi)$, et puisque $\Pi$ et $I - \Pi$ sont des projections
complémentaires,
\[
A^n = 2^n\,\Pi + (-1)^n (I - \Pi),
\qquad\text{c'est-à-dire}\qquad
(A^n)_{ij} = \frac{2^n}3 + (-1)^n\Bigl(\delta_{ij} -
\frac13\Bigr),
\]
ce qui donne les deux formules affichées. En $n = 2$~: diagonale
$(4 + 2)/3 = 2$ (chemins $i \to \ell \to i$ pour les deux voisins
$\ell$)~; hors diagonale $(4 - 1)/3 = 1$ (l'unique chemin $i \to
\ell \to j$ par le troisième sommet).

\textbf{17.} Soit $w_n^{(0)}, w_n^{(1)}$ le nombre de mots admissibles
de longueur $n$ finissant par $0$, resp.\ par $1$. En ajoutant une
lettre~: un $0$ peut suivre n'importe quoi, un $1$ seulement un $0$~:
\[
\begin{pmatrix} w_{n+1}^{(0)}\\ w_{n+1}^{(1)}\end{pmatrix}
= \begin{pmatrix} 1 & 1\\ 1 & 0\end{pmatrix}
\begin{pmatrix} w_n^{(0)}\\ w_n^{(1)}\end{pmatrix}.
\]
En sommant, $w_{n+2} = w_{n+1} + w_n$ (ou~: conditionner sur la
première lettre). Avec $w_1 = 2$, $w_2 = 3$~: $w_n = F_{n+2}$ par
récurrence ($F_3 = 2$, $F_4 = 3$, même récurrence). Croissance~: les
racines de $X^2 - X - 1$ sont $\varphi > \abs\psi$
(\cref{exo:b2:reduction:5}), et la composante en $\varphi$ est non
nulle (les $w_n$ sont positifs et $\psi^n \to 0$), donc la question 11
donne $w_{n+1}/w_n \to \varphi = \frac{1 + \sqrt5}2$.

\textbf{18.} $A = \begin{psmallmatrix} 0&1&0\\ 1&0&1\\
0&1&0\end{psmallmatrix}$. Vérification~:
\[
A(1, \pm\sqrt2, 1)^{\mathsf T}
= (\pm\sqrt2, 2, \pm\sqrt2)^{\mathsf T}
= \pm\sqrt2\,(1, \pm\sqrt2, 1)^{\mathsf T},
\qquad
A(1, 0, -1)^{\mathsf T} = 0 :
\]
\omterm{def:b2:reduction:eigen}{valeurs propres} $\sqrt2, -\sqrt2, 0$ ($= 2\cos\frac\pi4,
2\cos\frac{3\pi}4, 2\cos\frac\pi2$). On décompose $e_1$ sur la base
propre et on lit la troisième coordonnée, ou l'on utilise la
symétrie~: avec $v_\pm = (1, \pm\sqrt2, 1)$, $v_0 = (1, 0, -1)$, on
vérifie $e_1 = \frac14 v_+ + \frac14 v_- + \frac12 v_0$, donc pour $n
\geq 1$
\[
(A^n)_{13} = \Bigl(\tfrac14(\sqrt2)^n v_+ +
\tfrac14(-\sqrt2)^n v_- + 0\Bigr)_{\!3}
= \frac{(\sqrt2)^n + (-\sqrt2)^n}{4},
\]
nul pour $n$ impair (graphe biparti~: les extrémités sont à distance
paire), et $2\cdot 2^{n/2}/4 = 2^{n/2 - 1}$ pour $n$ pair. En $n =
4$~: $2^{1} = 2$, en accord avec les deux chemins $1\,2\,1\,2\,3$ et
$1\,2\,3\,2\,3$.

\textbf{19.} Les chemins fermés de longueur $n$ issus de $i$ sont
$(A^n)_{ii}$~; en sommant sur $i$ on obtient $\operatorname{tr}(A^n)$.
En trigonalisant $A$ (sur $\C$), $A^n$ est triangulaire de diagonale
$\lambda_i^n$~: $\operatorname{tr}(A^n) = \sum_i\lambda_i^n$.
Triangle~: $\operatorname{tr}(A^n) = 3\,\frac{2^n + 2(-1)^n}3 =
2^n + 2(-1)^n = 2^n + (-1)^n + (-1)^n$~: le \omterm{def:b2:reduction:eigen}{spectre} $\{2, -1, -1\}$,
cohérent avec la question 16.

\textbf{20.} Les colonnes de $W^{\mathsf T}$~: $W^{\mathsf T}e_i =
e_{i-1}$ pour $i \geq 1$ et $W^{\mathsf T}e_0 = e_{k-1}$~; en
réétiquetant dans l'\omterm{def:b2:structures:generated}{ordre} $e_0, e_1, \dots$ c'est exactement la
matrice compagnon de $X^k - 1$ ($a_0 = 1$, autres $a_m = 0$).
Question 2~: $\chi_{W} = \chi_{W^{\mathsf T}} = X^k - 1 = \mu_{W}$.
Les racines $\omega^j$ ($j = 0, \dots, k-1$) sont les $k$ racines
$k$-ièmes distinctes de l'unité~: $W$ est \omterm{def:b2:reduction:diag}{diagonalisable} (question 3,
ou \cref{exo:b2:reduction:8}~: $W^k = I$). Vecteurs propres~: $Wf_j
= \sum_m \omega^{-jm}e_{m+1} = \sum_{m'}\omega^{-j(m'-1)}e_{m'}
= \omega^j f_j$.

\textbf{21.} Les circulantes sont des polynômes en $W$, et les
polynômes en une matrice fixée commutent entre eux. Chaque $f_j$ est
un vecteur propre de toute puissance~: $W^m f_j = \omega^{jm}f_j$,
donc
\[
Cf_j = \sum_m c_m\omega^{jm} f_j = \widehat c(\omega^j)\,f_j :
\]
la base $(f_0, \dots, f_{k-1})$ (libre~: Vandermonde en les
$\omega^{-j}$ distincts, \cref{exo:b2:linalg:11}) diagonalise toute
circulante d'un coup, avec les \omterm{def:b2:reduction:eigen}{valeurs propres} annoncées.

\textbf{22.} Le \omterm{def:b2:linalg:det}{déterminant} est le produit des \omterm{def:b2:reduction:eigen}{valeurs propres}
(diagonaliser)~: $\det C = \prod_{j}\widehat c(\omega^j)$. Pour $k =
3$, $c_0 = a$, $c_1 = b$, $c_2 = c$ et $\omega = j =
\eu^{2\iu\pi/3}$~:
\[
\det C = (a + b + c)(a + bj + cj^2)(a + bj^2 + cj^4),
\]
et $j^4 = j$~: exactement la factorisation de
l'\cref{exo:b2:linalg:8}.

\textbf{23.} $M = \frac12(W + W^{-1})$ est une circulante ($W^{-1} =
W^{k-1}$), de \omterm{def:b2:reduction:eigen}{valeurs propres} $\frac12(\omega^j + \omega^{-j}) =
\cos\frac{2\pi j}k$ sur la même base $f_j$. Coordonnées~: écrivons
$x^{(0)} = \sum_j \alpha_j f_j$. Les coordonnées de $f_j$ somment à
$\sum_m \omega^{-jm}$, qui vaut $k$ pour $j = 0$ et $0$ sinon (somme
géométrique de raison $\omega^{-j} \neq 1$). En sommant les
coordonnées de $x^{(0)}$~: $\sum_m x^{(0)}_m = \alpha_0\,k$, donc
$\alpha_0 = \frac1k\sum_m x^{(0)}_m$, la moyenne.

\textbf{24.} $x^{(n)} = M^nx^{(0)} = \sum_j
\alpha_j\cos^n\bigl(\tfrac{2\pi j}k\bigr)f_j$. Pour $k$ impair,
$\abs{\cos(2\pi j/k)} < 1$ pour tout $j \neq 0$ (l'angle n'est jamais
$0$ ni $\pi$), donc tous les termes sauf $j = 0$ tendent vers $0$~:
$x^{(n)} \to \alpha_0 f_0$, le vecteur constant égal à la moyenne ---
moyenner sur un anneau impair égalise. Pour $k = 4$ les \omterm{def:b2:reduction:eigen}{valeurs
propres} sont $1, 0, -1, 0$~: le terme $j = 2$ $\alpha_2(-1)^nf_2$ avec
$f_2 = (1, -1, 1, -1)^{\mathsf T}$ oscille indéfiniment. L'obstruction
est la moyenne \omterm{def:b2:linalg:alternating}{alternée}~: en multipliant les coordonnées de $x^{(0)}$
par $(-1)^m$ et en sommant, le même calcul de somme géométrique donne
$\sum_m (-1)^mx^{(0)}_m = 4\alpha_2$~: le processus converge si et
seulement si $x^{(0)}_0 - x^{(0)}_1 + x^{(0)}_2 - x^{(0)}_3 = 0$, et
converge alors vers la moyenne.

\textbf{25.} La matrice compagnon convertit une récurrence scalaire
d'\omterm{def:b2:structures:generated}{ordre} $k$ en une récurrence vectorielle d'\omterm{def:b2:structures:generated}{ordre} un, de sorte que les
formules closes deviennent des énoncés sur $C^n$ --- le terrain de la
réduction (questions 1--5). Le lemme de décomposition des noyaux est
pure \omterm{def:b2:structures:algebra}{algèbre} polynomiale (Bézout plus commutation), donc il scinde
$\ker P(S)$ même si $\mathcal{S}$ est de dimension infinie
(questions 7--9). Les \omterm{def:b2:reduction:eigen}{valeurs propres} dominantes gouvernent la
croissance parce que toute autre contribution est géométriquement
négligeable après normalisation --- ce qui explique aussi pourquoi
l'erreur de Pell décroît au carré de la racine dominante
(questions 11--14). Les puissances de la matrice d'adjacence comptent
les chemins parce que la multiplication matricielle somme sur les
sommets intermédiaires, donc les \omterm{def:b2:reduction:eigen}{spectres} comptent les chemins fermés
(questions 15--19). Les matrices qui commutent partagent une base
propre, et une base de Fourier diagonalise alors toute l'\omterm{def:b2:structures:algebra}{algèbre} des
circulantes d'un seul geste (questions 20--24). Sommets~: le théorème
fondamental des récurrences linéaires (question 9)~; et pour les
matrices positives, la raison pour laquelle les racines dominantes
comme $\varphi$ ou $1 + \sqrt2$ sont automatiquement réelles,
positives et simples est le théorème de Perron--Frobenius, prouvé dans
le volume de l'année 3.
\end{solution}
